\section{Results and Discussion}
\label{sec:results}

\subsection{Macroscopic response}
\label{subsec:macro_response}

% Effective stress paths and stress-strain for uni, single-8, double-8
% Liquefaction resistance curves (CSR vs N_L)
% Key result: uni > single-8 > double-8 in N_L
% Strain path development: single-8 vs double-8 differences

\subsection{Energy interpretation}
\label{subsec:energy}

% Cumulative shear work comparison
% Similar total energy at liquefaction despite different N_L
% Implication: directionality affects rate of energy input, not total energy threshold

\subsection{Microscopic fabric interpretation}
\label{subsec:micro_fabric}

\paragraph{Coordination number evolution.}
% Zm degradation rates differ: double-8 fastest, uni slowest
% Consistent with N_L ordering

\paragraph{Fabric anisotropy evolution.}
% Fc evolution: double-8 develops highest anisotropy
% Physical interpretation: directional changes promote contact alignment

\paragraph{Contact density distribution.}
% Full directional evolution for three loading paths
% Morphological differences during loading, convergence near liquefaction

\subsection{Discussion}
\label{subsec:discussion}

% Why does directional change accelerate liquefaction?
% Connection between loading history and fabric degradation rate
% Comparison with previous DEM studies (Wei et al., Yang et al.)
% Implications for seismic hazard assessment
